\chapter{Test di ipotesi: il test di Neyman-Pearson}

La regola MAP del capitolo precedente ha bisogno di due ingredienti: i costi delle decisioni sbagliate e le probabilità
a priori delle ipotesi. Ci sono però molte situazioni in cui bisogna scegliere fra due ipotesi senza avere né l'uno né
l'altro. Le slide ne elencano alcune:
\begin{itemize}
  \item la rivelazione precoce delle minacce alla sicurezza di un'area pattugliata;
  \item la rivelazione di un'intrusione in server o domini protetti in Internet;
  \item la rivelazione (e la localizzazione) di un ostacolo nei sistemi avanzati di assistenza alla guida;
  \item il controllo del traffico aereo;
  \item un innumerevole insieme di applicazioni militari.
\end{itemize}
In tutti questi casi è praticamente impossibile assegnare un costo a un'interpretazione sbagliata dello «stato della
natura», cioè a una decisione errata fra le due ipotesi «tutto normale» e «sta succedendo qualcosa». E ha poco senso
assegnare delle probabilità a priori quando nei dati possono comparire delle «anomalie statistiche» mai viste prima.
Serve un criterio diverso: è quello dei \textbf{test di ipotesi}.

% ============================================================
\section{Ipotesi nulla, falsi allarmi e mancate rivelazioni}
% ============================================================

\dfn{Ipotesi nulla e ipotesi alternativa}{
  L'\term{ipotesi nulla}, tradizionalmente indicata con $H_0$, afferma che i dati osservati $\vett{x}$ sono la
  realizzazione di un vettore aleatorio con una distribuzione condizionata nota, con pmf $p_{\vett{X} \mid
  H_0}(\vett{x} \mid H_0)$ o PDF $f_{\vett{X} \mid H_0}(\vett{x} \mid H_0)$. Il \term{test di ipotesi} decide se, dati i
  dati osservati, l'ipotesi nulla deve essere rifiutata in favore di una legge diversa, l'\term{ipotesi alternativa}
  $H_1$, con pmf $p_{\vett{X} \mid H_1}(\vett{x} \mid H_1)$ o PDF $f_{\vett{X} \mid H_1}(\vett{x} \mid H_1)$.
}

L'ipotesi nulla è lo stato «normale», quello che si dà per buono finché i dati non lo smentiscono. Come nella
classificazione binaria bisogna dividere i dati in due regioni di decisione, $\Omega_0$ (si accetta $H_0$) e $\Omega_1$
(si rifiuta $H_0$ e si decide $H_1$), ma la struttura bayesiana non si può più usare, per mancanza di informazione a
priori. Si guarda invece separatamente ai due modi di sbagliare.

\dfn{Errori di prima e di seconda specie}{
  \begin{itemize}
    \item L'\term{errore di prima specie} (\emph{type I}) consiste nel rifiutare $H_0$ quando $H_0$ è vera. La sua
          probabilità, detta \term{probabilità di falso allarme} e indicata con $\alpha$, è
          \[ \alpha = P\{D(\vett{X}) = 1 \mid H_0\} = \begin{cases} \displaystyle\int_{\Omega_1} f_{\vett{X} \mid
             H_0}(\vett{x} \mid H_0) \dd\vett{x} & \text{dati continui,} \\[3mm] \displaystyle\sum_{\vett{x} \in
             \Omega_1} p_{\vett{X} \mid H_0}(\vett{x} \mid H_0) & \text{dati discreti.} \end{cases} \]
    \item L'\term{errore di seconda specie} (\emph{type II}) consiste nel rifiutare $H_1$ quando $H_1$ è vera. La sua
          probabilità, detta \term{probabilità di mancata rivelazione}, si indica con $\beta$. Il suo complemento a uno,
          detto \term{potenza del test} (o probabilità di corretta rivelazione), è
          \[ 1 - \beta = P\{D(\vett{X}) = 1 \mid H_1\} = \begin{cases} \displaystyle\int_{\Omega_1} f_{\vett{X} \mid
             H_1}(\vett{x} \mid H_1) \dd\vett{x} & \text{dati continui,} \\[3mm] \displaystyle\sum_{\vett{x} \in
             \Omega_1} p_{\vett{X} \mid H_1}(\vett{x} \mid H_1) & \text{dati discreti.} \end{cases} \]
  \end{itemize}
}

I quattro casi possibili si riassumono in una tabella. Si pensi a un allarme antifurto: $H_0$ è «nessun ladro», $H_1$ è
«c'è un ladro».
\begin{center}
\renewcommand{\arraystretch}{1.5}
\begin{tabular}{@{}l|c|c@{}}
  & \textbf{è vera $H_0$} (nessun ladro) & \textbf{è vera $H_1$} (c'è un ladro) \\ \hline
  \textbf{si decide $H_0$} (silenzio) & \cellcolor{mspiC!15} decisione giusta, prob. $1 - \alpha$ &
  \cellcolor{mspiB!20} mancata rivelazione (II specie), prob. $\beta$ \\ \hline
  \textbf{si decide $H_1$} (allarme) & \cellcolor{mspiB!20} falso allarme (I specie), prob. $\alpha$ &
  \cellcolor{mspiC!15} rivelazione, prob. $1 - \beta$ (potenza)
\end{tabular}
\end{center}
Le due colonne sono due mondi separati: in ciascuna le probabilità sommano a 1. Non si possono rendere piccoli entrambi
gli errori a piacere: allargando la regione $\Omega_1$ (un allarme più sensibile) la potenza cresce, ma crescono anche
i falsi allarmi; restringendola succede il contrario.

% ============================================================
\section{Il test di Neyman-Pearson}
% ============================================================

\dfn{Test di Neyman-Pearson}{
  Un \term{test di Neyman-Pearson} è il risultato della seguente ottimizzazione vincolata: si determina $\Omega_1$ in
  modo che
  \[ \begin{cases} 1 - \beta \ (\text{potenza del test}) : & \text{massimizzata,} \\ \text{vincolo:} & \text{probabilità
     di errore di prima specie} \leq \alpha . \end{cases} \]
}

Si fissa dunque quanti falsi allarmi si è disposti a tollerare (per esempio $\alpha = 0{,}01$: un falso allarme ogni
cento prove in condizioni normali) e, fra tutti i test che rispettano questo vincolo, si sceglie il più potente.
L'esistenza della soluzione è la sostanza del lemma di Neyman-Pearson, che dice anche quale forma ha.

\teorema[neymanpearson]{Lemma di Neyman-Pearson}{
  Il test più potente fra quelli con probabilità di falso allarme non superiore ad $\alpha$ è basato sul rapporto di
  verosimiglianza:
  \[ L(\vett{x}) \decidi{H_1}{H_0} \eta , \qquad\text{dove}\qquad L(\vett{x}) = \begin{cases}
     \dfrac{f_{\vett{X} \mid H_1}(\vett{x} \mid H_1)}{f_{\vett{X} \mid H_0}(\vett{x} \mid H_0)} & \text{dati continui,}
     \\[7mm] \dfrac{p_{\vett{X} \mid H_1}(\vett{x} \mid H_1)}{p_{\vett{X} \mid H_0}(\vett{x} \mid H_0)} & \text{dati
     discreti,} \end{cases} \]
  e la soglia $\eta$ si sceglie come soluzione dell'equazione $P\{L(\vett{X}) > \eta \mid H_0\} = \alpha$.
}

\begin{dimostrazione}[nel caso continuo; nel discreto gli integrali diventano somme]
\passo{La funzione indicatrice} Per una regione $\Omega$ si indica con $\mathbf{1}_{\Omega}(\vett{x})$ la funzione che
vale 1 se $\vett{x} \in \Omega$ e 0 altrimenti (è la variabile indicatrice del Capitolo~18, vista come funzione dei
dati). Così $\int \mathbf{1}_{\Omega}\, f_{\vett{X} \mid H_i} \dd\vett{x}$ è la probabilità della regione $\Omega$ sotto
$H_i$. Scriviamo in breve $f_0$ e $f_1$ per le due densità.
\passo{I due test a confronto} Sia $\Omega^\star = \{\vett{x} : f_1(\vett{x}) > \eta f_0(\vett{x})\}$ la regione del
test di Neyman-Pearson, con falso allarme esattamente $\alpha$, e sia $\Omega$ la regione di un test qualsiasi con
falso allarme $\alpha_\Omega \leq \alpha$.
\passo{Un integrale che non può essere negativo} Si consideri
\[ I = \int \left[ \mathbf{1}_{\Omega^\star}(\vett{x}) - \mathbf{1}_{\Omega}(\vett{x}) \right] \left[ f_1(\vett{x}) -
   \eta f_0(\vett{x}) \right] \dd\vett{x} . \]
Dove la prima parentesi è positiva ($\vett{x}$ sta in $\Omega^\star$ ma non in $\Omega$), anche la seconda lo è, per
come è definita $\Omega^\star$; dove la prima è negativa ($\vett{x}$ sta in $\Omega$ ma non in $\Omega^\star$), la
seconda è negativa o nulla. Il prodotto non è mai negativo, quindi $I \geq 0$.
\passo{Si svolge l'integrale} Moltiplicando, ogni pezzo è una probabilità:
\[ I = \underbrace{\left[ (1 - \beta_{\Omega^\star}) - (1 - \beta_{\Omega}) \right]}_{\text{differenza di potenza}} -
   \eta \underbrace{\left[ \alpha - \alpha_{\Omega} \right]}_{\geq 0} \geq 0 . \]
\passo{Conclusione} Quindi la differenza di potenza è almeno $\eta (\alpha - \alpha_\Omega)$, che non è negativo
($\eta \geq 0$ e $\alpha_\Omega \leq \alpha$): nessun test che rispetti il vincolo è più potente di quello di
Neyman-Pearson.
\end{dimostrazione}

Il test ha la stessa forma della regola MAP del capitolo precedente: si confronta il rapporto di verosimiglianza con una
soglia. Cambia il modo di scegliere la soglia: nella regola MAP è fissata dalle probabilità a priori, $\eta =
\frac{P(H_0)}{P(H_1)}$; qui è fissata dal livello di falsi allarmi $\alpha$ che si è disposti a tollerare. Con dati
discreti può non esistere una soglia che dia esattamente $\alpha$ (la probabilità sale a gradini); si sceglie allora la
soglia che dà il più grande falso allarme non superiore ad $\alpha$.

\warning{Le trasformazioni crescenti non cambiano il test}{
  Se si applica una qualsiasi funzione strettamente crescente a entrambi i membri della disuguaglianza, l'ottimalità del
  test non cambia: le stesse $\vett{x}$ finiscono nella stessa regione. Per esempio si può usare il rapporto di
  verosimiglianza logaritmico $\Lambda(\vett{x}) = \ln L(\vett{x})$ e confrontarlo con la soglia $\ln \eta$: con le
  densità gaussiane il logaritmo trasforma gli esponenziali in polinomi e i conti si semplificano molto.
}

% ============================================================
\section{Esempio: test sulla media di una popolazione gaussiana}
% ============================================================

L'ipotesi nulla è che le osservazioni siano variabili gaussiane i.i.d. con media nulla e varianza $\sigma^2$,
$X_i \sim \mathcal{N}(0, \sigma^2)$ (solo rumore); l'ipotesi alternativa è $X_i \sim \mathcal{N}(\mu, \sigma^2)$ con
$\mu > 0$ (un segnale di ampiezza $\mu$ immerso nel rumore). È il problema di un radar che deve capire se c'è un
bersaglio.

\teorema[testnp]{Il test di Neyman-Pearson sulla media}{
  Con un livello di falso allarme $\alpha$ assegnato, il test di Neyman-Pearson è
  \[ \frac{1}{n} \sum_{k=1}^{n} x_k \decidi{H_1}{H_0} \eta' = Q^{-1}(\alpha)\, \frac{\sigma}{\sqrt{n}} , \]
  e la sua potenza è
  \[ 1 - \beta = Q\left( Q^{-1}(\alpha) - \mu\, \frac{\sqrt{n}}{\sigma} \right) \;\xrightarrow{\;n \to \infty\;}\; 1 . \]
}

Qui $Q^{-1}$ è la \term{funzione inversa} della funzione $Q$: siccome $Q$ è strettamente decrescente (Capitolo~15), per
ogni $\alpha \in (0, 1)$ c'è un solo numero $x$ con $Q(x) = \alpha$, e lo si chiama $Q^{-1}(\alpha)$. Si legge sulla
tabella dei valori di $Q$ «al contrario», cercando $\alpha$ fra i risultati:
\begin{center}
\renewcommand{\arraystretch}{1.25}
\begin{tabular}{@{}c|cccc@{}}
  $\alpha$ & $0{,}1$ & $0{,}05$ & $0{,}01$ & $0{,}001$ \\ \hline
  $Q^{-1}(\alpha)$ & $1{,}282$ & $1{,}645$ & $2{,}326$ & $3{,}090$
\end{tabular}
\end{center}

\begin{dimostrazione}
\passo{Il rapporto di verosimiglianza} Come nel capitolo precedente, le costanti si semplificano:
\[ L(\vett{x}) = \frac{\prod_{k=1}^{n} e^{-\frac{(x_k - \mu)^2}{2\sigma^2}}}{\prod_{k=1}^{n} e^{-\frac{x_k^2}{2\sigma^2}}}
   = e^{\frac{1}{2\sigma^2} \sum_{k=1}^{n} \left[ x_k^2 - (x_k - \mu)^2 \right]} = e^{\frac{1}{2\sigma^2}
   \sum_{k=1}^{n} \left[ 2\mu x_k - \mu^2 \right]} , \]
perché $x_k^2 - (x_k - \mu)^2 = x_k^2 - x_k^2 - \mu^2 + 2\mu x_k$.
\passo{Il logaritmo} Prendendo il logaritmo di entrambi i lati e assorbendo in una nuova soglia $\eta'' = \ln \eta$
tutto ciò che non dipende dai dati:
\[ \Lambda(\vett{x}) = \frac{1}{2\sigma^2} \left[ 2\mu \sum_{k=1}^{n} x_k - n\mu^2 \right] \decidi{H_1}{H_0} \eta''
   \quad\iff\quad \mu \sum_{k=1}^{n} x_k \decidi{H_1}{H_0} \eta'' \sigma^2 + \frac{n \mu^2}{2} . \]
(Nelle slide 44--45 l'ipotesi sotto il simbolo è scritta $H_2$: va letta $H_0$.)
\passo{La media campione} Dividendo per $\mu n > 0$ il verso non cambia: $\frac{1}{n} \sum_k x_k \decidi{H_1}{H_0}
\eta'$, con $\eta' = \frac{1}{\mu n} \left( \eta'' \sigma^2 + \frac{n \mu^2}{2} \right)$. Invece di calcolare $\eta'$
da $\eta$, conviene sceglierlo direttamente in modo che il falso allarme valga $\alpha$.
\passo{La soglia} Sotto $H_0$, $Z_n = \frac{1}{n} \sum_k X_k$ è gaussiana con media nulla e varianza
$\frac{\sigma^2}{n}$ (Capitolo~18). Standardizzando,
\[ P\{Z_n > \eta' \mid H_0\} = Q\left( \frac{\eta'}{\sigma / \sqrt{n}} \right) = \alpha \quad\iff\quad
   \frac{\eta'}{\sigma / \sqrt{n}} = Q^{-1}(\alpha) \quad\iff\quad \eta' = Q^{-1}(\alpha)\, \frac{\sigma}{\sqrt{n}} . \]
\passo{La potenza} Sotto $H_1$, $Z_n$ è gaussiana con media $\mu$ e varianza $\frac{\sigma^2}{n}$, quindi
\[ 1 - \beta = P\{Z_n > \eta' \mid H_1\} = Q\left( \frac{\eta' - \mu}{\sigma / \sqrt{n}} \right) = Q\left(
   \eta' \frac{\sqrt{n}}{\sigma} - \mu \frac{\sqrt{n}}{\sigma} \right) = Q\left( Q^{-1}(\alpha) - \mu
   \frac{\sqrt{n}}{\sigma} \right) . \]
\passo{Il limite} Comunque si fissi un valore di $\alpha$ piccolo ma non nullo, $Q^{-1}(\alpha)$ è un numero grande ma
finito; per $n \to \infty$ l'argomento di $Q$ tende a $-\infty$ e $Q(-\infty) = 1$.
\end{dimostrazione}

\begin{figure}[H]
\centering
\begin{tikzpicture}
\begin{axis}[quaderno, axis y line=none, width=14cm, height=5.4cm, xmin=-1.1, xmax=1.65, ymin=0, ymax=1.9,
  xtick={0,0.411,0.5}, xticklabels={$0$,$\eta'$,$\mu$}, xlabel=$z$, legend style={font=\scriptsize, at={(0.99,0.99)},
  anchor=north east}]
  \addplot[fill=mspiB!45, draw=none, domain=0.411:1.6, samples=50, forget plot] {4*exp(-8*x^2)/sqrt(2*pi)} \closedcycle;
  \addplot[fill=mspiA!25, draw=none, domain=-1.1:0.411, samples=50, forget plot] {4*exp(-8*(x-0.5)^2)/sqrt(2*pi)}
    \closedcycle;
  \addplot[mspiB, domain=-1.1:1.6, samples=161] {4*exp(-8*x^2)/sqrt(2*pi)};
  \addplot[mspiA, domain=-1.1:1.6, samples=161] {4*exp(-8*(x-0.5)^2)/sqrt(2*pi)};
  \draw[black, dashed, line width=0.9pt] (axis cs:0.411,0) -- (axis cs:0.411,1.75) node[above, font=\scriptsize]
    {soglia};
  \draw[-{Stealth}] (axis cs:0.95,0.45) node[right, font=\scriptsize] {$\alpha = 0{,}05$ (falso allarme)} --
    (axis cs:0.55,0.08);
  \draw[-{Stealth}] (axis cs:-0.45,0.75) node[left, font=\scriptsize] {$\beta \approx 0{,}36$ (mancata rivelazione)}
    -- (axis cs:0.2,0.35);
  \legend{{$f_{Z_n \mid H_0}$: solo rumore}, {$f_{Z_n \mid H_1}$: segnale più rumore}}
\end{axis}
\end{tikzpicture}
\caption{Il test di Neyman-Pearson con $\sigma = 1$, $\mu = 0{,}5$, $n = 16$ e $\alpha = 0{,}05$: la soglia $\eta' =
1{,}645 \cdot \frac14 \approx 0{,}41$ lascia a destra il 5\% della campana di $H_0$. La potenza è l'area della campana di
$H_1$ a destra della soglia, $1 - \beta = Q(1{,}645 - 2) \approx 0{,}64$.}
\end{figure}

\begin{esempio}{Quante misure servono per un radar affidabile?}
Con $\sigma = 1$ e $\mu = 0{,}5$ (il segnale è metà del rumore) e $\alpha = 0{,}05$, la potenza vale
\[ 1 - \beta = Q\left( 1{,}645 - 0{,}5\sqrt{n} \right) : \]
con $n = 16$ misure $Q(-0{,}355) \approx 0{,}64$, con $n = 36$ $Q(-1{,}355) \approx 0{,}91$, con $n = 64$
$Q(-2{,}355) \approx 0{,}99$. Pretendere meno falsi allarmi costa potenza: con $\alpha = 0{,}01$ la soglia sale a
$Q^{-1}(0{,}01) = 2{,}326$ e, a parità di $n$, la potenza scende (curva più in basso nella figura).
\begin{center}
\begin{tikzpicture}
\begin{axis}[width=12cm, height=5.2cm, xmin=0, xmax=100, ymin=0, ymax=1.05, xtick={0,16,36,64,100},
  ytick={0,0.25,0.5,0.75,1}, yticklabels={$0$,$0{,}25$,$0{,}5$,$0{,}75$,$1$}, xlabel={numero di misure $n$},
  ylabel={potenza $1 - \beta$}, axis lines=left, tick label style={font=\scriptsize}, label style={font=\small},
  grid=major, grid style={mspigray!20}, legend style={font=\scriptsize, at={(0.98,0.03)}, anchor=south east}]
  \addplot[mspiA, line width=1.1pt] coordinates {(1,0.1261) (2,0.1742) (3,0.2180) (4,0.2595) (5,0.2992) (6,0.3372) (7,0.3737) (8,0.4088) (9,0.4424) (10,0.4746) (11,0.5054) (12,0.5347) (13,0.5627) (14,0.5894) (15,0.6147) (16,0.6388) (17,0.6616) (18,0.6831) (19,0.7035) (20,0.7228) (21,0.7410) (22,0.7581) (23,0.7743) (24,0.7895) (25,0.8038) (26,0.8172) (27,0.8298) (28,0.8416) (29,0.8526) (30,0.8630) (31,0.8727) (32,0.8817) (33,0.8902) (34,0.8981) (35,0.9054) (36,0.9123) (37,0.9187) (38,0.9247) (39,0.9302) (40,0.9354) (41,0.9402) (42,0.9447) (43,0.9489) (44,0.9527) (45,0.9563) (46,0.9596) (47,0.9627) (48,0.9656) (49,0.9682) (50,0.9707) (51,0.9729) (52,0.9750) (53,0.9770) (54,0.9788) (55,0.9805) (56,0.9820) (57,0.9834) (58,0.9847) (59,0.9859) (60,0.9871) (61,0.9881) (62,0.9891) (63,0.9899) (64,0.9907) (65,0.9915) (66,0.9922) (67,0.9928) (68,0.9934) (69,0.9939) (70,0.9944) (71,0.9949) (72,0.9953) (73,0.9957) (74,0.9960) (75,0.9964) (76,0.9967) (77,0.9970) (78,0.9972) (79,0.9974) (80,0.9977) (81,0.9978) (82,0.9980) (83,0.9982) (84,0.9983) (85,0.9985) (86,0.9986) (87,0.9987) (88,0.9988) (89,0.9989) (90,0.9990) (91,0.9991) (92,0.9992) (93,0.9993) (94,0.9993) (95,0.9994) (96,0.9994) (97,0.9995) (98,0.9995) (99,0.9996) (100,0.9996)};
  \addplot[mspiB, line width=1.1pt] coordinates {(1,0.0339) (2,0.0527) (3,0.0721) (4,0.0924) (5,0.1135) (6,0.1353) (7,0.1578) (8,0.1808) (9,0.2043) (10,0.2281) (11,0.2521) (12,0.2762) (13,0.3003) (14,0.3244) (15,0.3483) (16,0.3721) (17,0.3956) (18,0.4188) (19,0.4416) (20,0.4640) (21,0.4860) (22,0.5075) (23,0.5285) (24,0.5490) (25,0.5689) (26,0.5883) (27,0.6071) (28,0.6253) (29,0.6429) (30,0.6599) (31,0.6764) (32,0.6922) (33,0.7074) (34,0.7221) (35,0.7362) (36,0.7497) (37,0.7627) (38,0.7751) (39,0.7870) (40,0.7984) (41,0.8093) (42,0.8196) (43,0.8295) (44,0.8390) (45,0.8480) (46,0.8565) (47,0.8647) (48,0.8724) (49,0.8797) (50,0.8867) (51,0.8933) (52,0.8996) (53,0.9055) (54,0.9112) (55,0.9165) (56,0.9215) (57,0.9263) (58,0.9308) (59,0.9350) (60,0.9390) (61,0.9428) (62,0.9464) (63,0.9497) (64,0.9529) (65,0.9559) (66,0.9587) (67,0.9613) (68,0.9638) (69,0.9661) (70,0.9683) (71,0.9704) (72,0.9723) (73,0.9742) (74,0.9759) (75,0.9775) (76,0.9790) (77,0.9804) (78,0.9817) (79,0.9829) (80,0.9841) (81,0.9851) (82,0.9861) (83,0.9871) (84,0.9880) (85,0.9888) (86,0.9896) (87,0.9903) (88,0.9910) (89,0.9916) (90,0.9922) (91,0.9927) (92,0.9932) (93,0.9937) (94,0.9942) (95,0.9946) (96,0.9950) (97,0.9953) (98,0.9956) (99,0.9960) (100,0.9962)};
  \legend{{$\alpha = 0{,}05$}, {$\alpha = 0{,}01$}}
\end{axis}
\end{tikzpicture}
\end{center}
\end{esempio}

\info{Una soglia che non dipende da $\mu$}{
  La soglia $\eta' = Q^{-1}(\alpha) \frac{\sigma}{\sqrt{n}}$ dipende solo dal rumore e dal livello di falso allarme, non
  dall'ampiezza $\mu$ del segnale. Quindi lo stesso test (confrontare la media delle misure con $\eta'$) è il migliore
  per \emph{qualunque} $\mu > 0$, anche se l'ampiezza del bersaglio non è nota. Da $\mu$ dipende solo la potenza:
  bersagli deboli si vedono peggio.
}

Le prestazioni di un test si riassumono nella \term{curva ROC} (\emph{Receiver Operating Characteristic}): al variare
della soglia, si disegna la potenza $1 - \beta$ in funzione del falso allarme $\alpha$. Per il test sulla media la
curva dipende solo da $d = \mu \frac{\sqrt{n}}{\sigma}$, la distanza fra le due campane misurata in deviazioni standard
della media campione: $1 - \beta = Q\left( Q^{-1}(\alpha) - d \right)$.

\begin{figure}[H]
\centering
\begin{tikzpicture}
\begin{axis}[width=8.4cm, height=8cm, xmin=0, xmax=1, ymin=0, ymax=1, xtick={0,0.25,0.5,0.75,1},
  ytick={0,0.25,0.5,0.75,1}, xticklabels={$0$,$0{,}25$,$0{,}5$,$0{,}75$,$1$},
  yticklabels={$0$,$0{,}25$,$0{,}5$,$0{,}75$,$1$}, xlabel={falso allarme $\alpha$}, ylabel={potenza $1 - \beta$},
  axis lines=left, tick label style={font=\scriptsize}, label style={font=\small}, grid=major,
  grid style={mspigray!20}, legend style={font=\scriptsize, at={(0.98,0.03)}, anchor=south east}]
  \addplot[mspiC, line width=1.1pt] coordinates {(1.0000,1.0000) (1.0000,1.0000) (1.0000,1.0000) (1.0000,1.0000) (1.0000,1.0000) (1.0000,1.0000) (1.0000,1.0000) (1.0000,1.0000) (1.0000,1.0000) (1.0000,1.0000) (1.0000,1.0000) (1.0000,1.0000) (1.0000,1.0000) (1.0000,1.0000) (1.0000,1.0000) (0.9999,1.0000) (0.9998,1.0000) (0.9997,1.0000) (0.9995,1.0000) (0.9992,1.0000) (0.9987,1.0000) (0.9978,1.0000) (0.9965,1.0000) (0.9946,1.0000) (0.9918,1.0000) (0.9878,1.0000) (0.9821,1.0000) (0.9744,1.0000) (0.9641,1.0000) (0.9505,1.0000) (0.9332,1.0000) (0.9115,1.0000) (0.8849,1.0000) (0.8531,1.0000) (0.8159,1.0000) (0.7734,0.9999) (0.7257,0.9998) (0.6736,0.9997) (0.6179,0.9995) (0.5596,0.9992) (0.5000,0.9987) (0.4404,0.9978) (0.3821,0.9965) (0.3264,0.9946) (0.2743,0.9918) (0.2266,0.9878) (0.1841,0.9821) (0.1469,0.9744) (0.1151,0.9641) (0.0885,0.9505) (0.0668,0.9332) (0.0495,0.9115) (0.0359,0.8849) (0.0256,0.8531) (0.0179,0.8159) (0.0122,0.7734) (0.0082,0.7257) (0.0054,0.6736) (0.0035,0.6179) (0.0022,0.5596) (0.0013,0.5000) (0.0008,0.4404) (0.0005,0.3821) (0.0003,0.3264) (0.0002,0.2743) (0.0001,0.2266) (0.0000,0.1841) (0.0000,0.1469) (0.0000,0.1151) (0.0000,0.0885) (0.0000,0.0668) (0.0000,0.0495) (0.0000,0.0359) (0.0000,0.0256) (0.0000,0.0179) (0.0000,0.0122) (0.0000,0.0082) (0.0000,0.0054) (0.0000,0.0035) (0.0000,0.0022) (0.0000,0.0013) (0.0000,0.0008) (0.0000,0.0005) (0.0000,0.0003) (0.0000,0.0002) (0.0000,0.0001) (0.0000,0.0000)};
  \addplot[mspiA, line width=1.1pt] coordinates {(1.0000,1.0000) (1.0000,1.0000) (1.0000,1.0000) (1.0000,1.0000) (1.0000,1.0000) (1.0000,1.0000) (1.0000,1.0000) (1.0000,1.0000) (1.0000,1.0000) (1.0000,1.0000) (1.0000,1.0000) (1.0000,1.0000) (1.0000,1.0000) (1.0000,1.0000) (1.0000,1.0000) (0.9999,1.0000) (0.9998,1.0000) (0.9997,1.0000) (0.9995,1.0000) (0.9992,1.0000) (0.9987,1.0000) (0.9978,1.0000) (0.9965,1.0000) (0.9946,1.0000) (0.9918,1.0000) (0.9878,1.0000) (0.9821,1.0000) (0.9744,1.0000) (0.9641,0.9999) (0.9505,0.9999) (0.9332,0.9998) (0.9115,0.9996) (0.8849,0.9993) (0.8531,0.9989) (0.8159,0.9981) (0.7734,0.9970) (0.7257,0.9953) (0.6736,0.9929) (0.6179,0.9893) (0.5596,0.9842) (0.5000,0.9772) (0.4404,0.9678) (0.3821,0.9554) (0.3264,0.9394) (0.2743,0.9192) (0.2266,0.8944) (0.1841,0.8643) (0.1469,0.8289) (0.1151,0.7881) (0.0885,0.7422) (0.0668,0.6915) (0.0495,0.6368) (0.0359,0.5793) (0.0256,0.5199) (0.0179,0.4602) (0.0122,0.4013) (0.0082,0.3446) (0.0054,0.2912) (0.0035,0.2420) (0.0022,0.1977) (0.0013,0.1587) (0.0008,0.1251) (0.0005,0.0968) (0.0003,0.0735) (0.0002,0.0548) (0.0001,0.0401) (0.0000,0.0287) (0.0000,0.0202) (0.0000,0.0139) (0.0000,0.0094) (0.0000,0.0062) (0.0000,0.0040) (0.0000,0.0026) (0.0000,0.0016) (0.0000,0.0010) (0.0000,0.0006) (0.0000,0.0003) (0.0000,0.0002) (0.0000,0.0001) (0.0000,0.0001) (0.0000,0.0000) (0.0000,0.0000) (0.0000,0.0000) (0.0000,0.0000) (0.0000,0.0000) (0.0000,0.0000) (0.0000,0.0000)};
  \addplot[mspiB, line width=1.1pt] coordinates {(1.0000,1.0000) (1.0000,1.0000) (1.0000,1.0000) (1.0000,1.0000) (1.0000,1.0000) (1.0000,1.0000) (1.0000,1.0000) (1.0000,1.0000) (1.0000,1.0000) (1.0000,1.0000) (1.0000,1.0000) (1.0000,1.0000) (1.0000,1.0000) (1.0000,1.0000) (1.0000,1.0000) (0.9999,1.0000) (0.9998,1.0000) (0.9997,1.0000) (0.9995,1.0000) (0.9992,1.0000) (0.9987,1.0000) (0.9978,0.9999) (0.9965,0.9999) (0.9946,0.9998) (0.9918,0.9997) (0.9878,0.9994) (0.9821,0.9990) (0.9744,0.9984) (0.9641,0.9974) (0.9505,0.9960) (0.9332,0.9938) (0.9115,0.9906) (0.8849,0.9861) (0.8531,0.9798) (0.8159,0.9713) (0.7734,0.9599) (0.7257,0.9452) (0.6736,0.9265) (0.6179,0.9032) (0.5596,0.8749) (0.5000,0.8413) (0.4404,0.8023) (0.3821,0.7580) (0.3264,0.7088) (0.2743,0.6554) (0.2266,0.5987) (0.1841,0.5398) (0.1469,0.4801) (0.1151,0.4207) (0.0885,0.3632) (0.0668,0.3085) (0.0495,0.2578) (0.0359,0.2119) (0.0256,0.1711) (0.0179,0.1357) (0.0122,0.1056) (0.0082,0.0808) (0.0054,0.0606) (0.0035,0.0446) (0.0022,0.0322) (0.0013,0.0228) (0.0008,0.0158) (0.0005,0.0107) (0.0003,0.0071) (0.0002,0.0047) (0.0001,0.0030) (0.0000,0.0019) (0.0000,0.0011) (0.0000,0.0007) (0.0000,0.0004) (0.0000,0.0002) (0.0000,0.0001) (0.0000,0.0001) (0.0000,0.0000) (0.0000,0.0000) (0.0000,0.0000) (0.0000,0.0000) (0.0000,0.0000) (0.0000,0.0000) (0.0000,0.0000) (0.0000,0.0000) (0.0000,0.0000) (0.0000,0.0000) (0.0000,0.0000) (0.0000,0.0000) (0.0000,0.0000) (0.0000,0.0000)};
  \addplot[mspiric, line width=1.1pt] coordinates {(1.0000,1.0000) (1.0000,1.0000) (1.0000,1.0000) (1.0000,1.0000) (1.0000,1.0000) (1.0000,1.0000) (1.0000,1.0000) (1.0000,1.0000) (1.0000,1.0000) (1.0000,1.0000) (1.0000,1.0000) (1.0000,1.0000) (1.0000,1.0000) (1.0000,1.0000) (1.0000,1.0000) (0.9999,1.0000) (0.9998,1.0000) (0.9997,1.0000) (0.9995,0.9999) (0.9992,0.9999) (0.9987,0.9998) (0.9978,0.9996) (0.9965,0.9993) (0.9946,0.9989) (0.9918,0.9981) (0.9878,0.9970) (0.9821,0.9953) (0.9744,0.9929) (0.9641,0.9893) (0.9505,0.9842) (0.9332,0.9772) (0.9115,0.9678) (0.8849,0.9554) (0.8531,0.9394) (0.8159,0.9192) (0.7734,0.8944) (0.7257,0.8643) (0.6736,0.8289) (0.6179,0.7881) (0.5596,0.7422) (0.5000,0.6915) (0.4404,0.6368) (0.3821,0.5793) (0.3264,0.5199) (0.2743,0.4602) (0.2266,0.4013) (0.1841,0.3446) (0.1469,0.2912) (0.1151,0.2420) (0.0885,0.1977) (0.0668,0.1587) (0.0495,0.1251) (0.0359,0.0968) (0.0256,0.0735) (0.0179,0.0548) (0.0122,0.0401) (0.0082,0.0287) (0.0054,0.0202) (0.0035,0.0139) (0.0022,0.0094) (0.0013,0.0062) (0.0008,0.0040) (0.0005,0.0026) (0.0003,0.0016) (0.0002,0.0010) (0.0001,0.0006) (0.0000,0.0003) (0.0000,0.0002) (0.0000,0.0001) (0.0000,0.0001) (0.0000,0.0000) (0.0000,0.0000) (0.0000,0.0000) (0.0000,0.0000) (0.0000,0.0000) (0.0000,0.0000) (0.0000,0.0000) (0.0000,0.0000) (0.0000,0.0000) (0.0000,0.0000) (0.0000,0.0000) (0.0000,0.0000) (0.0000,0.0000) (0.0000,0.0000) (0.0000,0.0000) (0.0000,0.0000) (0.0000,0.0000)};
  \addplot[black, dashed, domain=0:1, samples=2] {x};
  \legend{{$d = 3$}, {$d = 2$}, {$d = 1$}, {$d = 0{,}5$}, {tirare a caso}}
\end{axis}
\end{tikzpicture}
\hspace{0.4cm}
\begin{minipage}[b]{6.2cm}\small
Ogni punto di una curva corrisponde a una soglia: abbassandola si sale lungo la curva, con più rivelazioni ma anche più
falsi allarmi. Più la curva si avvicina all'angolo in alto a sinistra ($\alpha = 0$, $1 - \beta = 1$), migliore è il
test. La diagonale tratteggiata è un test che decide tirando una moneta truccata: rivela con la stessa probabilità con
cui dà falsi allarmi. Aumentando il numero di misure $n$ cresce $d$, e la curva sale.
\vspace{0.6cm}
\end{minipage}
\caption{Le curve ROC del test sulla media gaussiana per diversi valori di $d = \mu\sqrt{n}/\sigma$.}
\end{figure}

\gbox{In sintesi}{mspiA!80!black}{
\begin{itemize}
  \item Senza costi e senza probabilità a priori si usa un \textbf{test di ipotesi}: ipotesi nulla $H_0$ contro ipotesi
        alternativa $H_1$.
  \item \textbf{Errore di prima specie} (falso allarme): si rifiuta $H_0$ vera, probabilità $\alpha$. \textbf{Errore di
        seconda specie} (mancata rivelazione): si rifiuta $H_1$ vera, probabilità $\beta$. \textbf{Potenza}: $1 -
        \beta$.
  \item \textbf{Neyman-Pearson}: si massimizza la potenza con il vincolo falso allarme $\leq \alpha$. Il test ottimo è
        $L(\vett{x}) \decidi{H_1}{H_0} \eta$, con $\eta$ tale che $P\{L(\vett{X}) > \eta \mid H_0\} = \alpha$; si può
        usare $\ln L$.
  \item Media gaussiana: $\frac1n \sum_k x_k \decidi{H_1}{H_0} Q^{-1}(\alpha) \frac{\sigma}{\sqrt{n}}$, con potenza
        $Q\left( Q^{-1}(\alpha) - \mu \frac{\sqrt{n}}{\sigma} \right) \to 1$. La curva ROC riassume il compromesso fra
        $\alpha$ e $1 - \beta$.
\end{itemize}
}
